% MAS1004 Assignment 1: long report.
%
% You fill this in with your agent, in your repository. AGENTS.md tells the
% agent how. Red text in the PDF is what still has to be replaced.
%
% - Pictures go in figures/. Copy them from results/, which is not in git.
% - Outputs go in outputs/. run.py saves what it prints in
%   results/<tag>_output.txt, and PLAN.md saves the other outputs in results/
%   with `| tee results/...`. Copy them here. \outputfile{outputs/check.txt}
%   prints a file exactly as it is. Every number in this report has to come
%   from such an output.
%
% To make the PDF: compress this folder into a zip, and on Overleaf choose
% New Project, Upload Project. Set Menu, Compiler to XeLaTeX, and Recompile.

\documentclass[11pt]{article}
\usepackage{mas1004}

\reporttitle{Assignment 1: Long report}
\studentname{\fillin{your name}}
\studentid{\fillin{your student id}}

\begin{document}
\makeheader

% ---------------------------------------------------------------------------
\section*{Problem 1. Categories and images}

\subsection*{Why these categories}
All three categories are shoes, and each search phrase names a tread pattern rather than a kind of
shoe: \texttt{converse\_diamond\_sole\_shape}, \texttt{nike\_air\_force\_circular\_sole\_shape} and
\texttt{vans\_honeycomb\_sole\_shape}. I chose them together because the object is the same in all
three, so the only thing left for a classifier to work with is the geometry of the sole: the
diamond chevron grid of a Converse, the circular pivot-point pattern of an Air Force~1, and the
honeycomb waffle grid of a Vans.
\fillin{I still owe two or three sentences here, in my own words, about why I care about these
three; this part of Problem 1 is mine alone.}

\subsection*{What I expected the model to use}
\fillin{Paste the guess I wrote down in Step 1, before any training. I cannot reconstruct it from
the repository, and the assignment wants the guess as it was written then.}

\subsection*{Images per category}
150 images came back for each of the three phrases, 450 in all. Run 1 read
\texttt{150 images, 3 classes} and one line per class, 150 each
(\texttt{results/run\_output.txt}); the searches themselves are in
\texttt{results/cleaning/temp\_logs/collect\_log.txt}.

% ---------------------------------------------------------------------------
\section*{Problem 2. Preparing an image}

\sidebyside{figures/prepare_original.jpg}{the image as downloaded}
           {figures/prepare_result.png}{what \texttt{prepare\_image} made of it}

The picture above is \texttt{data/clean/nike\_air\_force\_circular\_sole\_shape/0009.jpg}, a
photograph of a Nike sole that is 1600 x 566 pixels, wide and short. \texttt{prepare\_image}
resizes its shorter side to 256, keeping the shape (723 x 256), and then cuts the middle
224 x 224 square out of it. 499 of the 723 pixels of the long side are thrown away and only
31.0\% of the width survives. So on a photo shaped like this one the crop keeps the middle of the
sole and cuts off its heel end and its toe end, which is the part of the picture where the
circular pattern is most visible. The numbers come from running
\texttt{torchvision.transforms.Resize} and \texttt{CenterCrop} with \texttt{RESIZE} and
\texttt{CROP} from \texttt{src/data.py}.

% ---------------------------------------------------------------------------
\section*{Problem 3. First training runs}

\begin{center}\footnotesize
\begin{tabular}{llllrrrr}
\toprule
tag & start & trained & epochs & lr & trainable parameters & train acc. & test acc. \\
\midrule
run & ImageNet & all layers & 10 & 0.0001 & 11,178,051 & 99.2\% & 85.6\% \\
scratch & random & all layers & 10 & 0.0001 & 11,178,051 & 74.7\% & 44.4\% \\
frozen & ImageNet & last layer & 10 & 0.001 & 1,539 & 88.3\% & 77.8\% \\
long & ImageNet & all layers & 20 & 0.0003 & 11,178,051 & 100.0\% & 84.4\% \\
\bottomrule
\end{tabular}
\end{center}

\picturefile{figures/run_curves.png}

The training loss falls from 0.92 to 0.02 in the first four epochs and then stays between 0.01 and
0.03 for the remaining six; the test loss falls from 0.56 to about 0.38 by epoch 5 and then stays
between 0.38 and 0.42. So the loss went down and then flattened out, and it flattened on the test
side: after epoch 5 the model is no longer learning anything that transfers to the test images, it
is only getting better at the 360 training images. (\texttt{results/run\_output.txt}, epochs 1-10.)

The gap is 99.2\% training against 85.6\% test, 13.6 points. The model reproduces the images it
was trained on almost exactly and still gets about one test image in seven wrong, so part of what
it learned is the 360 training images themselves rather than the pattern that separates them.

\begin{center}\picturefile[0.5\linewidth]{figures/run_confusion.png}\end{center}

The two it mixes up most are nike and vans, in both directions: of its 30 nike test images it gets
24 right and 4 of the misses it calls vans, and of the 30 vans it gets 24 right, calling 3 of the
misses nike and 3 converse. Converse is almost perfect (29 of 30, one called vans). That is what I
expected rather than a surprise: the circular Air Force pattern and the honeycomb Vans pattern are
both dense, round, textured rubber, while the diamond chevron of a Converse is the odd one out.
(The confusion matrix of every run, printed by the same code path as
\texttt{evaluate.confusion\_matrix} on each run's own test split, is in
\texttt{results/confusions2.txt}.)

From ImageNet weights the model reaches 85.6\%; from random numbers it reaches 44.4\%, barely above
the 33.3\% you get by guessing one of three. That 41.2-point gap is not a training-budget artefact:
the random-weights run gets the same 10 epochs, the same learning rate and the same 360 images.
ImageNet weights bring edges, textures and shapes learned from over a million photographs, and my
three classes are separated by exactly that kind of low-level texture, which 360 images cannot
teach from nothing.

Training only the last layer changes 1,539 of the 11,178,051 parameters, 0.014\% of them, and
reaches 77.8\% against 85.6\% for the whole network. So it recovers most of the accuracy for a
fraction of the work, but the last 7.8 points need the earlier layers to move as well.

\subsection*{Output of the runs}
\outputfile{outputs/run_output.txt}
\outputfile{outputs/scratch_output.txt}
\outputfile{outputs/frozen_output.txt}
\outputfile{outputs/long_output.txt}

% ---------------------------------------------------------------------------
\section*{Problem 4. Cleaning}

\subsection*{My rule}
An image belongs in its category only if it is an ordinary photograph in which the shoe's sole is
the subject and the tread geometry (the diamond, circular or honeycomb pattern that names the
class) is visible, in focus and large enough to crop a usable square from. I remove everything
else: pictures of the whole shoe with the sole not visible, tread that is a different pattern or
from a shoe that is not this one, images where the pattern comes from something that is not a
shoe, drawings, logos, text or collage graphics, and duplicates of an image I have already kept.
I wrote it down in \texttt{results/cleaning/rule.txt} before running \texttt{clean.py suspects},
and I wrote it before removing anything.

\subsection*{Three kinds of junk}
\begin{enumerate}
  \item A whole shoe with the sole out of frame (\texttt{no-tread}).
        \texttt{converse\_diamond\_sole\_shape/0036.jpg} is a tidy product photo of a grey high-top
        seen from the side. The photo is fine; there is simply no tread in it, so it cannot teach
        the diamond pattern. This was by far the most common kind: 186 of the 246 removals.
  \item Not a photograph at all (\texttt{not-a-photo}).
        \texttt{nike\_air\_force\_circular\_sole\_shape/0072.jpg} is a technical drawing of a shoe
        with measurement lines around it. It shows the shape of the sole clearly, which is exactly
        why the first draft of my rule let it through, so the rule now says photograph.
  \item Not the object (\texttt{wrong-object}).
        \texttt{vans\_honeycomb\_sole\_shape/0062.jpg} is a yellow bus. Searching "honeycomb sole"
        also returned honeycomb flooring and tiles, a box truck, and listing screenshots of vans,
        which is why the rule now says the pattern has to come from a shoe.
\end{enumerate}
\picturefile[0.31\linewidth]{figures/junk_no_tread.jpg}\hfill
\picturefile[0.31\linewidth]{figures/junk_not_a_photo.jpg}\hfill
\picturefile[0.31\linewidth]{figures/junk_wrong_object.jpg}

\subsection*{What I removed}
\outputfile{outputs/clean_count.txt}

\texttt{clean.py suspects} found 22 clusters of near copies, covering 59 of the 450 images, from
114 pairs above the 0.95 similarity at which it calls two images a near copy. One image of each
cluster had to go, so 37 of those 59 came out: 22 of them with the reason \texttt{duplicate} and
the other 15 because they had already gone for a stronger reason (\texttt{no-tread} or
\texttt{not-a-photo}). The pairs are drawn in \texttt{figures/copies.png}. The biggest cluster was
16 nearly identical photographs of the same brown sole in vans (0001, 0002, 0003, 0005, 0007, 0009,
0010, 0015-0018 and 0021-0025): I kept 0003 and removed the other 15. The clusters come from a
union-find over the \texttt{copy of} rows of
\texttt{results/cleaning/suspects/suspects.csv}.

For each class the tool also listed the images whose own label it found least likely. I removed 12
of the first 20 suspects in converse, 13 in nike and 7 in vans; all 60 are listed with the reason
each one got in \texttt{results/cleaning/first20.txt}. Two things are worth saying about that
number. The suspects did include junk I had missed: they find the 16-copy brown-sole cluster,
which I had read as sixteen ordinary sole photographs in my own pass, and nine of the vans top 20
sit on the two look-sheets I did not manage to review by eye. I left those nine in rather than
remove images I had never looked at, so they are the one part of the cleaning that is unfinished;
they are the last block in \texttt{results/cleaning/first20.txt}. And a suspect is not junk: the
converse images I kept from its top 20 are real photographs of real soles at unusual angles (a
pair of soles held in a hand, four soles lined up on a table, a white sole with a star on it),
which the network finds unlike the rest of the class even though my rule keeps them. That is the
point the assignment makes, and it is why the tool never removed anything by itself.

\subsection*{Before and after cleaning}
\outputfile{outputs/compare.txt}

Test accuracy went from 85.6\% (\texttt{run}, 90 test images, before cleaning) to 80.0\%
(\texttt{clean}, 40 test images, after): it went \emph{down} by 5.6 points, and that is a real
result rather than a mistake in the numbers. The test sets are not the same images. Cleaning took
246 of the 450 images out, and the test set lost its share of them, but what it lost was the easy
part: a photograph of a yellow bus, a CAD drawing or a listing screenshot is nowhere near the
other two classes, so an old test set full of them scores high partly for the wrong reason and
unfairly flatters the old model. Cleaning also leaves less to learn from: the clean model trains
on 164 images instead of 360, and its training loss is already 0.03 by epoch 5, with the test
accuracy no longer moving after epoch 2. So the honest reading is that the 85.6\% was measured on
a test set containing junk the model could not fail on, and the 80.0\% is measured on almost
nothing but real sole photographs. Which of the two models is actually better is decided on the
new images, which cleaning does not touch: \texttt{check.py --compare run clean}, printed above.

The clean model's own confusion matrix, in \texttt{results/confusions2.txt}, shows where its 8 test
errors sit: 5 of them are converse images, the class cleaning cut hardest (150 down to 62, so only
12 test images), and it gets 7 of those 12 right, 2 wrong as nike and 3 as vans. Vans is now
perfect (15 of 15). A class with 12 test images is a noisy measurement, so I am careful not to read
too much into 58\%, but it does match the fact that converse is the class where my rule removed the
largest share of images.

% ---------------------------------------------------------------------------
\section*{Problem 5. Images from a new source}

\fillin{Where your new images came from, and why you are sure none of them are
among your downloads.}

\outputfile{outputs/overlap.txt}

\fillin{The accuracy on your downloaded test set and on your new images, side
by side.}

\fillin{Five mistakes it was confident about, with the pictures. For each one,
what in the image you think pushed it the wrong way.}

\fillin{What is different about your new images: background, lighting, angle,
distance, what else is in the frame.}

% ---------------------------------------------------------------------------
\section*{Problem 6. The demo}

Address: \url{https://isabel-mg.github.io/mas1004-assignment1/} and the code is at
\url{https://github.com/Isabel-mg/mas1004-assignment1}. The model runs inside the page, so there is
no server and no photo leaves the machine. The badge is green: the self-test that the export wrote
reports that the page and Python agree with a largest gap of 0.0000, and
\texttt{docs/model.onnx} is 44.8 MB with 11,178,051 parameters.

The page has no "none of these" answer. It always returns the single class with the highest
probability, whatever you show it, so a picture of something else comes back as whichever of the
three sole patterns is nearest.
\fillin{Write here what you actually saw when you showed the live page something that is none of
the three categories, and the confidence it printed; this has to be your own test, on your phone.}

% ---------------------------------------------------------------------------
\section*{Problem 7. What the model looks at}

% Copy this subsection once for every guess.
\subsection*{Guess 1}

\fillin{The guess, and what in your images or your model's mistakes made you
think of it.}

\fillin{The change you chose, and why it tests the guess.}

My request to the agent, word for word:
\begin{quote}
\fillin{the request}
\end{quote}

% One \sidebyside for every image you changed.
\sidebyside{figures/guess1_original_1.jpg}{\fillin{original: the model's answer}}
           {figures/guess1_changed_1.jpg}{\fillin{changed: the model's answer}}

\fillin{What the results say about the guess.}

\subsection*{What I now think the model uses}
\fillin{What you now think the model uses to tell your categories apart, and
how sure you are.}

% ---------------------------------------------------------------------------
\section*{Appendix. Output of check.py}
\outputfile{outputs/check.txt}

\end{document}
